\documentclass[literature.tex]{subfiles}

\begin{document}
\drama{Revizor}{Nikolaj Vasiljevič Gogol}
\teorieA{
drama, satirická komedie (komedie omylů s groteskními prvky)
}{
Kritický realismus (zakladatel v ruské literatuře). Navazuje na romantismus, ale přechází k realistickému obrazu společnosti s prvky grotesky a satiry.
}{
Hovorový, živý jazyk s vulgarismy, archaismy a typickými ruskými příjmeními (mluvící jména – Skvoznik-Dmuchanovskij, Chlestakov). Dialogy odhalují charaktery – hejtman mluví důstojně, Chlestakov nafoukaně a nesouvisle.
}{
Ironie, opakování, kontrast (zdánlivá důstojnost vs. korupce), groteska.
}{
Satira, alegorie (městečko = celé Rusko), symbolismus (revizor = svědomí / Soudný den), hyperbola.
}{
Pět jednání s chronologickým postupem. Rychlé dialogy, scénické poznámky, groteskní situace. Smích z absurdity a lidské malosti („Nevrč, brachu, na zrcadlo, když máš křivou hubu“).
}{
HEJTMAN: Pánové, přijel revizor! Inkognito! A ještě s tajným příkazem!\\
BOBČINSKIJ a DOBČINSKIJ (současně): Ano, ano! Viděli jsme ho v hospodě!\\
HEJTMAN: Co teď? Musíme všechno uvést do pořádku... rychle!
}
\teorieB{
\wtem{Ivan Alexandrovič Chlestakov}{lehkomyslný, přitroublý petrohradský úředníček, povaleč a lhář. Využívá omylu a užívá si roli „revizora“.}
\wtem{Anton Antonovič (Hejtman)}{korupčník, lstivý, ale zbabělý. Bojí se kontroly, podplácí a lichotí.}
\wtem{Anna Andrejevna a Marie Antonovna}{hejtmanova žena a dcera -- marnivé, koketní, snící o Petrohradu.}
\wtem{Bobčinskij a Dobčinskij}{zvědaví, hloupí zemani, kteří celý omyl odstartují.}
\wtem{Osip}{Chlestakův sluha, chytřejší než pán, pragmatický.}
}{
V malém ruském městečku se rozšíří zpráva, že přijede tajný revizor. Zkorumpovaní úředníci (hejtman, soudce, školní inspektor, poštmistr aj.) se v panice snaží zakrýt nepořádek a úplatky. Dva zvědaví zemani považují za revizora lehkomyslného petrohradského úředníčka Chlestakova, který v hospodě nemá na zaplacení. Hejtman ho pozve k sobě domů, všichni mu lichotí a půjčují peníze. Chlestakov se vžije do role, naparuje se, dvoří se hejtmanově ženě i dceři a slibuje sňatek. Nakonec odjede. Poštmistr otevře jeho dopis a všichni zjistí, že to byl podvod. Závěrem přijíždí pravý revizor.
}{
Pět jednání, chronologický postup (děj proběhne během cca 24 hodin). Prostor: malé provinční ruské městečko (neurčené). Čas: 30. léta 19. století (vláda cara Mikuláše I.).
}{
Ostrá satira na korupci, byrokracii, maloměšťáctví a lidskou hloupost v carském Rusku. Motto „Nevrč na zrcadlo, když máš křivou hubu“ – kritika celé společnosti. Gogol ukazuje, jak strach z moci vede k podlézavosti a pokrytectví. Alegoricky: městečko = celé Rusko, Chlestakov = falešná autorita.
}
\historie{
30. léta 19. století – vláda cara Mikuláše I. (autokratický režim, cenzura, policejní dohled). Gogol kritizuje provinční korupci a byrokracii.
}{}{}{}{
Přechod od romantismu ke kritickému realismu. Námět údajně pochází od Puškina. Gogol ovlivnil Dostojevského, Čechova a celou ruskou literaturu.
}{
(1809–1852). Ukrajinského původu, ruský prozaik a dramatik. Studoval v Nižynu, žil v Petrohradě a Římě. Trápila ho deprese a náboženské krize. Zemřel v 42 letech po dlouhém půstu.
}{
Mezi nejznámější díla patří \textsc{Mrtvé duše}, \textsc{Večery na samotě u Dikanky}, \textsc{Taras Bulba}, povídky jako \textsc{Plášť} nebo \textsc{Nos}.
}\newpage
\kritika{}{}{
Nadčasová kritika korupce, byrokracie, podlézavosti a strachu z moci. Platí i dnes – „revizor“ může přijít kdykoli.
}
\nazor{
Revizor je geniální satirická komedie – směješ se, ale mrazí tě z toho, jak přesně Gogol vystihl lidskou malost. Chlestakov jako nafoukaný povaleč, který využije situaci, a hejtman, který se z paniky promění v podlézavce, jsou dokonalé typy. Nejsilnější je závěr s pravým revizorem – zrcadlo, které nikdo nechce vidět. Gogol dokázal kritizovat celou společnost vtipně a bez kázání. Rozhodně patří k vrcholům ruské dramatiky – čte se skvěle i dnes.
}
\explain{
\textbl{Satirická komedie --} zesměšňuje společenské nedostatky (korupce, byrokracie);
\textbl{Groteska --} přehánění a zkreslení reality pro komický efekt;
\textbl{Kritický realismus --} věrný, kritický obraz společnosti bez idealizace.
}
\wpage
\end{document}